% preamble/theorems.tex
% ============================================
% 定理环境体系 —— 墨绿 + 暗金
% 五层差异化 tcolorbox + 证明/解环境
% ============================================

% ---- 定理样式 ----
\declaretheoremstyle[
  spaceabove=6pt, spacebelow=6pt,
  headfont=\normalfont\bfseries,
  headpunct=,
  postheadspace=0.5em,
  bodyfont=\normalfont,
  headformat=\NAME
]{mycolored}

% ---- 声明定理类型（共享编号计数器）----
\declaretheorem[style=mycolored, name=, within=section]{theorem}
\declaretheorem[style=mycolored, name=, sibling=theorem]{lemma}
\declaretheorem[style=mycolored, name=, sibling=theorem]{corollary}
\declaretheorem[style=mycolored, name=, sibling=theorem]{proposition}
\declaretheorem[style=mycolored, name=, sibling=theorem]{definition}
\declaretheorem[style=mycolored, name=, sibling=theorem]{remark}
\declaretheorem[style=mycolored, name=, sibling=theorem]{example}
\declaretheorem[style=mycolored, name=, sibling=theorem]{construction}
\declaretheorem[style=mycolored, name=, sibling=theorem]{problem}
\declaretheorem[style=mycolored, name=, sibling=theorem]{axiom}
\declaretheorem[style=mycolored, name=, sibling=theorem]{prp}
\declaretheorem[style=mycolored, name=, sibling=theorem]{note}
\declaretheorem[style=mycolored, name=, sibling=theorem]{ill}

% Hyperref 锚点必须包含节号，否则定理计数器按节重置后会产生重名目标。
\renewcommand{\theHtheorem}{\theHsection.\arabic{theorem}}
\renewcommand{\theHlemma}{\theHsection.\arabic{lemma}}
\renewcommand{\theHcorollary}{\theHsection.\arabic{corollary}}
\renewcommand{\theHproposition}{\theHsection.\arabic{proposition}}
\renewcommand{\theHdefinition}{\theHsection.\arabic{definition}}
\renewcommand{\theHremark}{\theHsection.\arabic{remark}}
\renewcommand{\theHexample}{\theHsection.\arabic{example}}
\renewcommand{\theHconstruction}{\theHsection.\arabic{construction}}
\renewcommand{\theHproblem}{\theHsection.\arabic{problem}}
\renewcommand{\theHaxiom}{\theHsection.\arabic{axiom}}
\renewcommand{\theHprp}{\theHsection.\arabic{prp}}
\renewcommand{\theHnote}{\theHsection.\arabic{note}}
\renewcommand{\theHill}{\theHsection.\arabic{ill}}

% ---- 预计算编号工具 ----
\newcommand{\NextTheoremNum}{\thesection.\number\numexpr\value{theorem}+1\relax}
\newcommand{\NextLemmaNum}{\thesection.\number\numexpr\value{lemma}+1\relax}
\newcommand{\NextCorNum}{\thesection.\number\numexpr\value{corollary}+1\relax}
\newcommand{\NextPropNum}{\thesection.\number\numexpr\value{proposition}+1\relax}
\newcommand{\NextDefNum}{\thesection.\number\numexpr\value{definition}+1\relax}
\newcommand{\NextRemNum}{\thesection.\number\numexpr\value{remark}+1\relax}
\newcommand{\NextExNum}{\thesection.\number\numexpr\value{example}+1\relax}
\newcommand{\NextConNum}{\thesection.\number\numexpr\value{construction}+1\relax}
\newcommand{\NextAxiomNum}{\thesection.\number\numexpr\value{axiom}+1\relax}
\newcommand{\NextPrpNum}{\thesection.\number\numexpr\value{prp}+1\relax}
\newcommand{\NextNoteNum}{\thesection.\number\numexpr\value{note}+1\relax}
\newcommand{\NextIllNum}{\thesection.\number\numexpr\value{ill}+1\relax}
\newcommand{\NextProblemNum}{\thesection.\number\numexpr\value{problem}+1\relax}

% ============================================================
%  类别 I — 墨绿 GreenMain (30,80,60)
%  定理 / 推论 / 命题 / 原理
%  白底 + 墨绿边框 + 粗左侧线
% ============================================================
\AtBeginEnvironment{theorem}{%
	\theoremFont
	\begin{tcolorbox}[enhanced,breakable,
		colback=white,colframe=LogicBlue,
		arc=4pt,boxrule=0.6pt,leftrule=3pt,
		title=定理 \NextTheoremNum]}
\AtEndEnvironment{theorem}{\end{tcolorbox}}

\AtBeginEnvironment{corollary}{%
	\theoremFont
	\begin{tcolorbox}[enhanced,breakable,
		colback=white,colframe=LogicBlue,
		arc=4pt,boxrule=0.6pt,leftrule=3pt,
		title=推论 \NextCorNum]}
\AtEndEnvironment{corollary}{\end{tcolorbox}}

\AtBeginEnvironment{proposition}{%
	\theoremFont
	\begin{tcolorbox}[enhanced,breakable,
		colback=white,colframe=LogicBlue,
		arc=4pt,boxrule=0.6pt,leftrule=3pt,
		title=命题 \NextPropNum]}
\AtEndEnvironment{proposition}{\end{tcolorbox}}

\AtBeginEnvironment{prp}{%
	\theoremFont
	\begin{tcolorbox}[enhanced,breakable,
		colback=white,colframe=LogicBlue,
		arc=4pt,boxrule=0.6pt,leftrule=3pt,
		title=原理 \NextPrpNum]}
\AtEndEnvironment{prp}{\end{tcolorbox}}

% ============================================================
%  类别 II — 深墨绿 GreenDark (20,55,40)
%  公理 / 问题 / 引理 / 注记
%  白底 + 深墨绿边框 + 粗左侧线
% ============================================================
\AtBeginEnvironment{axiom}{%
	\theoremFont
	\begin{tcolorbox}[enhanced,breakable,
		colback=white,colframe=LogicDarkBlue,
		arc=4pt,boxrule=0.6pt,leftrule=3pt,
		title=公理 \NextAxiomNum]}
\AtEndEnvironment{axiom}{\end{tcolorbox}}

\AtBeginEnvironment{problem}{%
	\theoremFont
	\begin{tcolorbox}[enhanced,breakable,
		colback=white,colframe=LogicDarkBlue,
		arc=4pt,boxrule=0.6pt,leftrule=3pt,
		title=问题 \NextProblemNum]}
\AtEndEnvironment{problem}{\end{tcolorbox}}

\AtBeginEnvironment{lemma}{%
	\theoremFont
	\begin{tcolorbox}[enhanced,breakable,
		colback=white,colframe=LogicDarkBlue,
		arc=4pt,boxrule=0.6pt,leftrule=3pt,
		title=引理 \NextLemmaNum]}
\AtEndEnvironment{lemma}{\end{tcolorbox}}

\AtBeginEnvironment{remark}{%
	\theoremFont
	\begin{tcolorbox}[enhanced,breakable,
		colback=white,colframe=LogicDarkBlue,
		arc=4pt,boxrule=0.6pt,leftrule=3pt,
		title=注记 \NextRemNum]}
\AtEndEnvironment{remark}{\end{tcolorbox}}

% ============================================================
%  类别 III — 暗金 GoldMain (175,145,60)
%  定义 —— 暗金边框 + 淡金底，突出基石概念
% ============================================================
\AtBeginEnvironment{definition}{%
	\theoremFont
	\begin{tcolorbox}[enhanced,breakable,
		colback=GoldPale!30,colframe=GoldMain,
		arc=4pt,boxrule=0.6pt,leftrule=3pt,
		title={$\lozenge$\enspace 定义 \NextDefNum}]}
\AtEndEnvironment{definition}{\end{tcolorbox}}

% ============================================================
%  类别 IV — 浅绿 GreenLight (120,160,135)
%  构造 / 记号 / 说明
%  白底 + 浅绿边框 + 细左侧线
% ============================================================
\AtBeginEnvironment{construction}{%
	\theoremFont
	\begin{tcolorbox}[enhanced,breakable,
		colback=white,colframe=LogicLightBlue,
		arc=4pt,boxrule=0.6pt,leftrule=3pt,
		title=构造 \NextConNum]}
\AtEndEnvironment{construction}{\end{tcolorbox}}

\AtBeginEnvironment{note}{%
	\theoremFont
	\begin{tcolorbox}[enhanced,breakable,
		colback=white,colframe=LogicLightBlue,
		arc=4pt,boxrule=0.6pt,leftrule=3pt,
		title=记号 \NextNoteNum]}
\AtEndEnvironment{note}{\end{tcolorbox}}

\AtBeginEnvironment{ill}{%
	\theoremFont
	\begin{tcolorbox}[enhanced,breakable,
		colback=white,colframe=LogicLightBlue,
		arc=4pt,boxrule=0.6pt,leftrule=3pt,
		title=说明 \NextIllNum]}
\AtEndEnvironment{ill}{\end{tcolorbox}}

% ============================================================
%  类别 V — 示例
%  无色框，保留“示例 + 编号”的普通标题
% ============================================================
\AtBeginEnvironment{example}{%
	\theoremFont
	\par\medskip
	\noindent\textbf{示例 \NextExNum}\par\smallskip}
\AtEndEnvironment{example}{\par\medskip}

% ============================================================
%  解 (Solution) 环境 —— 独立计数器，楷体
% ============================================================
\newcounter{solution}
\renewcommand{\thesolution}{\arabic{solution}}
\newenvironment{solution}{%
	\par\proofFont
	\refstepcounter{solution}%
	\noindent\textbf{解.}\quad
}{%
	\hfill $\square$\par
}

% ============================================================
%  证明 (Proof) 环境 —— 无色块，楷体
% ============================================================
\renewenvironment{proof}[1][证明]{%
	\par\proofFont
	\noindent\textbf{#1.}\quad
}{%
	\hfill $\square$\par
}

% ---- 星号版本（按章编号）----
\declaretheorem[numberwithin=chapter, name=定理, sibling=theorem]{theorem*}
\declaretheorem[numberwithin=chapter, name=问题, sibling=problem]{problem*}

% 行距
\setstretch{1.2}
